\section{Generating the Single-Qutrit Clifford Group}
\label{sec:clifford-gen}

Both the physical RB and the logical LRB-D protocols sample uniformly from the single-qutrit Clifford group, so we need an explicit catalog of its elements, a realization of each element as a circuit, and a fast way to look up inverses for the recovery gate. This appendix records how that catalog is constructed.

Modulo global phases, the single-qutrit Clifford group $\mathcal{C}_3$ contains $216$ elements. It factors into the qutrit Pauli group and the quotient $\mathcal{C}_3 / \mathcal{P}_3$, which is isomorphic to the symplectic group $Sp(2,\mathbb{Z}_3) = SL(2,\mathbb{Z}_3)$ of order $24$ \cite{Appleby2005,Hostens2005}. Counting phase-free Pauli representatives $X^a Z^b$ with $a,b \in \mathbb{Z}_3$ gives $9$ elements, and $9 \times 24 = 216$. We exploit this factorization and build the group in two stages, first the $24$ coset representatives and then their products with the $9$ Pauli monomials.

\paragraph{Coset representatives from a Cayley graph.}
The quotient $\mathcal{C}_3/\mathcal{P}_3$ is generated by the images of $H$ and $S$. We represent each generator by its conjugation action on the eight non-identity Pauli classes $\{X, Z, XZ, X^2, Z^2, X^2Z, Z^2X, X^2Z^2\}$, which realizes the quotient faithfully as a permutation group on eight symbols. We then perform a breadth-first traversal of the Cayley graph of this group. Starting from the identity, every known element is multiplied on the right by each generator, and the first time a new element is produced we record its label as the label of its parent extended by the corresponding generator symbol. Because the traversal is breadth-first, each of the $24$ elements receives a minimal-length word in $H$ and $S$, and the traversal terminates once all $24$ are found. Repeated symbols are then compressed into power notation, so for instance $SSHS$ becomes $S^2HS$. The graph explored by this traversal is shown in \autoref{fig:cayley-c3modp3}, with every node carrying the canonical word that the traversal assigned to it.

\begin{figure}[H]
  \centering
  \resizebox{0.92\textwidth}{!}{%
\begin{tikzpicture}[>={Stealth[length=2.2mm,width=1.6mm]},
  cnode/.style={circle,draw=black,fill=white,inner sep=1.1pt,minimum size=7.5mm,font=\scriptsize},
  Hedge/.style={->,HCol,line width=0.75pt},
  Sedge/.style={->,SwapCol,line width=0.75pt}]
  \node[cnode] (n0) at (2.83,5.09) {$I$};
  \node[cnode] (n1) at (-1.32,6.5) {$H$};
  \node[cnode] (n2) at (2.41,0.84) {$S$};
  \node[cnode] (n3) at (1.74,6.4) {$H^2$};
  \node[cnode] (n4) at (-0.84,2.41) {$HS$};
  \node[cnode] (n5) at (-2.41,-0.84) {$SH$};
  \node[cnode] (n6) at (6.5,1.32) {$S^2$};
  \node[cnode] (n7) at (-2.49,5.26) {$H^3$};
  \node[cnode] (n8) at (1.0,2.35) {$H^2S$};
  \node[cnode] (n9) at (2.35,-1.0) {$HSH$};
  \node[cnode] (n10) at (-5.09,2.83) {$HS^2$};
  \node[cnode] (n11) at (-6.5,-1.32) {$SHS$};
  \node[cnode] (n12) at (5.09,-2.83) {$S^2H$};
  \node[cnode] (n13) at (-2.35,1.0) {$H^3S$};
  \node[cnode] (n14) at (-1.0,-2.35) {$H^2SH$};
  \node[cnode] (n15) at (5.26,2.49) {$H^2S^2$};
  \node[cnode] (n16) at (6.4,-1.74) {$HSHS$};
  \node[cnode] (n17) at (-6.4,1.74) {$SHSH$};
  \node[cnode] (n18) at (-2.83,-5.09) {$SHS^2$};
  \node[cnode] (n19) at (1.32,-6.5) {$S^2HS$};
  \node[cnode] (n20) at (0.84,-2.41) {$H^3SH$};
  \node[cnode] (n21) at (-5.26,-2.49) {$H^2SHS$};
  \node[cnode] (n22) at (2.49,-5.26) {$HSHS^2$};
  \node[cnode] (n23) at (-1.74,-6.4) {$S^2HSH$};
  \draw[Hedge] (n0) to[bend left=16] (n1);
  \draw[Sedge] (n0) to[bend left=16] (n2);
  \draw[Hedge] (n1) to[bend left=10] (n3);
  \draw[Sedge] (n1) to[bend left=16] (n4);
  \draw[Hedge] (n2) to[bend left=16] (n5);
  \draw[Sedge] (n2) to[bend left=16] (n6);
  \draw[Hedge] (n3) to[bend left=16] (n7);
  \draw[Sedge] (n3) to[bend left=16] (n8);
  \draw[Hedge] (n4) to[bend left=16] (n9);
  \draw[Sedge] (n4) to[bend left=16] (n10);
  \draw[Hedge] (n5) to[bend left=16] (n8);
  \draw[Sedge] (n5) to[bend left=16] (n11);
  \draw[Hedge] (n6) to[bend left=16] (n12);
  \draw[Sedge] (n6) to[bend left=16] (n0);
  \draw[Hedge] (n7) to[bend left=16] (n0);
  \draw[Sedge] (n7) to[bend left=16] (n13);
  \draw[Hedge] (n8) to[bend left=16] (n14);
  \draw[Sedge] (n8) to[bend left=16] (n15);
  \draw[Hedge] (n9) to[bend left=16] (n13);
  \draw[Sedge] (n9) to[bend left=16] (n16);
  \draw[Hedge] (n10) to[bend left=16] (n11);
  \draw[Sedge] (n10) to[bend left=16] (n1);
  \draw[Hedge] (n11) to[bend left=10] (n17);
  \draw[Sedge] (n11) to[bend left=16] (n18);
  \draw[Hedge] (n12) to[bend left=16] (n15);
  \draw[Sedge] (n12) to[bend left=16] (n19);
  \draw[Hedge] (n13) to[bend left=16] (n20);
  \draw[Sedge] (n13) to[bend left=16] (n17);
  \draw[Hedge] (n14) to[bend left=16] (n2);
  \draw[Sedge] (n14) to[bend left=16] (n21);
  \draw[Hedge] (n15) to[bend left=16] (n16);
  \draw[Sedge] (n15) to[bend left=16] (n3);
  \draw[Hedge] (n16) to[bend left=10] (n6);
  \draw[Sedge] (n16) to[bend left=16] (n22);
  \draw[Hedge] (n17) to[bend left=16] (n21);
  \draw[Sedge] (n17) to[bend left=16] (n7);
  \draw[Hedge] (n18) to[bend left=16] (n19);
  \draw[Sedge] (n18) to[bend left=16] (n5);
  \draw[Hedge] (n19) to[bend left=10] (n23);
  \draw[Sedge] (n19) to[bend left=16] (n20);
  \draw[Hedge] (n20) to[bend left=16] (n4);
  \draw[Sedge] (n20) to[bend left=16] (n12);
  \draw[Hedge] (n21) to[bend left=16] (n10);
  \draw[Sedge] (n21) to[bend left=16] (n23);
  \draw[Hedge] (n22) to[bend left=16] (n18);
  \draw[Sedge] (n22) to[bend left=16] (n9);
  \draw[Hedge] (n23) to[bend left=16] (n22);
  \draw[Sedge] (n23) to[bend left=16] (n14);
\end{tikzpicture}
  }
  \caption{Cayley graph of the quotient $\mathcal{C}_3/\mathcal{P}_3 \cong SL(2,\mathbb{Z}_3)$ generated by the traversal described in the text. Each of the $24$ nodes is one coset representative, labeled by the power-compressed word assigned to it when the breadth-first traversal first reached it. Orange arrows denote right multiplication by $H$ and purple arrows denote right multiplication by $S$. The $S$ arrows close into eight directed $3$-cycles and the $H$ arrows close into six directed $4$-cycles, reflecting the orders of the two generators, and every element is reachable from the identity $I$ within five generator applications.}
  \label{fig:cayley-c3modp3}
\end{figure}


\paragraph{Full expressions.}
The full catalog is the set of products $P \cdot C$, where $P$ ranges over the nine Pauli monomials and $C$ over the $24$ coset words. After normalizing the identity cases and removing duplicates, this yields exactly $216$ canonical expression strings, such as $X^2Z^2S^2HSH$. These strings are the sampling alphabet for every benchmark in this work. For the physical RB baseline, each string is parsed into its symbol and power tokens and the corresponding primitive gates are multiplied out into a single unencoded qutrit gate.

\paragraph{Logical realization.}
For the encoded protocol, each primitive symbol in an expression is replaced by its transversal logical implementation on the folded $[[5,1,2]]_3$ patch. The logical $\bar{H}$ and $\bar{S}$ are the constant-depth circuits of \autoref{fig: logical H} and \autoref{fig: logical S}, and the logical Paulis act as $\bar{X} = X_0 \otimes X_3$ and $\bar{Z} = Z_0 \otimes Z_2$. Since the expression strings are written as operator products, with the leftmost factor acting last, the substitution walks each string in reverse so that the resulting circuit applies the factors in the correct temporal order. Finally, runs of consecutive single-qutrit gates on the same physical wire are merged into single composite gates between the entangling operations. This keeps the depth of each logical Clifford close to the depth of its generating set rather than growing with the word length.

\paragraph{Inverses.}
The recovery gate $\bar{C}_{inv} = (\prod_{i=1}^m \bar{C}_i)^{-1}$ is never synthesized at run time. Instead, the group multiplication over the catalog is precomputed, and each catalog element is stored together with the expression of its inverse. Constructing the inverse of a sampled sequence then reduces to a table lookup followed by the same logical realization as any other group element.

The construction is implemented once in Cirq for the unencoded and encoded gate objects, and the resulting expression catalog together with its per-symbol gate decompositions is what the \texttt{sdim} simulation pipeline of \autoref{sub:sim_framework} consumes.
